\chapter{随机过程}

\section{随机过程基础}

\subsection{定义与分类}

随机过程是定义在概率空间上的时间索引随机变量族 $\{X_t\}_{t \in T}$。按状态空间和时间索引分类：

\begin{itemize}[leftmargin=*]
  \item 离散时间、离散状态：马尔可夫链（如信用评级迁移）
  \item 离散时间、连续状态：GARCH、ARMA
  \item 连续时间、连续状态：布朗运动、扩散过程（金融衍生品定价的核心）
  \item 连续时间、离散状态：泊松过程（跳跃风险建模）
\end{itemize}

\subsection{马尔可夫性质}

若过程的未来仅依赖于当前状态而不依赖历史，则为马尔可夫过程：

\begin{equation}
  P(X_{t+s} \in A \mid \mathcal{F}_t) = P(X_{t+s} \in A \mid X_t)
\end{equation}

许多金融模型基于此假设：资产价格的未来分布仅依赖于当前价格。

\subsection{鞅 (Martingale)}

鞅是公平博弈的数学形式化：

\begin{equation}
  \mathbb{E}[X_{t+s} \mid \mathcal{F}_t] = X_t
\end{equation}

在风险中性测度下，折现资产价格为鞅——这是现代衍生品定价的数学基石。

\section{布朗运动}

\subsection{标准布朗运动}

布朗运动（Wiener Process）$W_t$ 的性质：

\begin{enumerate}[leftmargin=*]
  \item $W_0 = 0$ a.s.
  \item \textbf{独立增量}：对 $0 \leq s < t$，$W_t - W_s$ 独立于 $\mathcal{F}_s$
  \item \textbf{平稳增量}：$W_t - W_s \sim \mathcal{N}(0, t-s)$
  \item \textbf{连续样本路径}：以概率 1，$t \mapsto W_t$ 连续
\end{enumerate}

\begin{keybox}
布朗运动的两个关键性质：(1) 处处连续但处处不可微；(2) 在有限区间上的二次变差为区间长度：$\langle W \rangle_t = t$。
\end{keybox}

\subsection{布朗运动的变换}

\begin{itemize}[leftmargin=*]
  \item \textbf{标度性质}：$cW_{t/c^2}$ 也是布朗运动
  \item \textbf{反射原理}：用于计算首达时分布
  \item \textbf{时间反演}：$t W_{1/t}$ 也是布朗运动
\end{itemize}

\section{伊藤积分}

\subsection{动机与定义}

伊藤积分解决了积分对象不可微的问题。对于简单过程：

\begin{equation}
  \int_0^T \phi_t dW_t = \sum_{i=0}^{n-1} \phi_{t_i}(W_{t_{i+1}} - W_{t_i})
\end{equation}

关键：积分值取左端点（非可料），这是伊藤积分区别于 Stratonovich 积分的核心。

\subsection{伊藤引理}

若 $X_t$ 满足 $dX_t = \mu_t dt + \sigma_t dW_t$，$f(t,x)$ 为 $C^{1,2}$ 函数，则：

\begin{equation}
  df(t,X_t) = \left(\frac{\partial f}{\partial t} + \mu_t\frac{\partial f}{\partial x} + \frac{1}{2}\sigma_t^2\frac{\partial^2 f}{\partial x^2}\right)dt + \sigma_t\frac{\partial f}{\partial x}dW_t
\end{equation}

$\frac{1}{2}\sigma_t^2\frac{\partial^2 f}{\partial x^2}$ 项是二阶修正——这是随机微积分与普通微积分的本质差异，来源于 $(dW_t)^2 = dt$。

\begin{example}[几何布朗运动的伊藤引理]
设 $S_t$ 遵循 $dS_t = \mu S_t dt + \sigma S_t dW_t$，令 $f(S) = \ln S$：
\begin{equation}
  d(\ln S_t) = \left(\mu - \frac{1}{2}\sigma^2\right)dt + \sigma dW_t
\end{equation}
因此 $\ln S_T \sim \mathcal{N}\left(\ln S_0 + (\mu - \sigma^2/2)T, \sigma^2 T\right)$。
\end{example}

\subsection{随机微分方程 (SDE)}

一般形式：
\begin{equation}
  dX_t = \mu(t,X_t)dt + \sigma(t,X_t)dW_t
\end{equation}

解的存在唯一性：Lipschitz 条件和线性增长条件（参考 Oksendal 的 SDE 理论）。

常见解析解：
\begin{itemize}[leftmargin=*]
  \item 几何布朗运动：$S_t = S_0 \exp((\mu - \sigma^2/2)t + \sigma W_t)$
  \item Ornstein-Uhlenbeck 过程：$X_t = X_0 e^{-\theta t} + \mu(1-e^{-\theta t}) + \sigma \int_0^t e^{-\theta(t-s)} dW_s$
\end{itemize}

\section{跳跃过程}

\subsection{泊松过程}

强度为 $\lambda$ 的泊松过程 $N_t$：
\begin{itemize}[leftmargin=*]
  \item $N_0 = 0$
  \item 独立增量
  \item $N_t - N_s \sim \text{Poisson}(\lambda(t-s))$
\end{itemize}

\subsection{复合泊松过程}

\begin{equation}
  J_t = \sum_{i=1}^{N_t} Y_i
\end{equation}

其中 $Y_i \sim \text{i.i.d.}$ 为跳跃幅度。

\subsection{跳跃扩散模型 (Merton, 1976)}

\begin{equation}
  \frac{dS_t}{S_{t-}} = \mu dt + \sigma dW_t + dJ_t
\end{equation}

其中 $J_t$ 为复合泊松过程，跳跃幅度通常取对数正态分布。Merton 模型可以生成隐含波动率微笑，弥补 BS 模型的局限。

\subsection{L\'evy 过程}

L\'evy 过程是具有平稳独立增量和随机连续的随机过程。L\'evy-Khintchine 表示：

\begin{equation}
  \mathbb{E}[e^{iuX_t}] = \exp\left[t\left(ibu - \frac{1}{2}\sigma^2 u^2 + \int_{\mathbb{R}}(e^{iux} - 1 - iux\mathbf{1}_{|x|<1})\nu(dx)\right)\right]
\end{equation}

其中 $(b,\sigma^2,\nu)$ 为特征三元组，$\nu$ 为 L\'evy 测度（控制跳跃频率和大小）。

常见金融 L\'evy 模型：VG（Variance Gamma）、NIG（Normal Inverse Gaussian）、CGMY。

\section{鞅表示与对冲}

\subsection{鞅表示定理}

在布朗运动过滤下，任意鞅 $M_t$ 可表示为：

\begin{equation}
  M_t = M_0 + \int_0^t H_s dW_s
\end{equation}

这保证了完备市场中所有未定权益可 Delta 对冲。

\subsection{市场完备性}

在一个存在 $d$ 个布朗运动和 $n$ 个风险资产的市场：
\begin{itemize}[leftmargin=*]
  \item $n \geq d$：市场完备（所有未定权益可复制）
  \item $n < d$：不完备市场（如带随机波动率的模型）
\end{itemize}

\section{测度变换与计价单位}

\subsection{概览}

测度变换（Measure Change）是 Q Quant 的核心操作：

\begin{theorem}[Girsanov]
若 $\frac{d\mathbb{Q}}{d\mathbb{P}} = \mathcal{E}\left(-\int_0^T \lambda_s dW_s\right)$，则：
\begin{equation}
  W_t^{\mathbb{Q}} = W_t^{\mathbb{P}} + \int_0^t \lambda_s ds
\end{equation}
在 $\mathbb{Q}$ 下是布朗运动。
\end{theorem}

在金融中，$\lambda$ 为夏普比率（Sharpe Ratio），即超额收益比波动率：

\begin{equation}
  \lambda_t = \frac{\mu_t - r}{\sigma_t}
\end{equation}

\subsection{数例}

$\mathbb{Q}$ 测度下股票 SDE：
\begin{equation}
  dS_t = r S_t dt + \sigma S_t dW_t^{\mathbb{Q}}
\end{equation}

漂移从 $\mu$ 变为 $r$（风险中性化）。Girsanov 定理从数学上证明：改变测度等价于改变漂移——波动率不变。

\section{Feynman-Kac 公式}

偏微分方程与随机过程的桥梁：

\begin{theorem}[Feynman-Kac]
若 $u(t,x)$ 满足：
\begin{equation}
  \frac{\partial u}{\partial t} + \mu(t,x)\frac{\partial u}{\partial x} + \frac{1}{2}\sigma^2(t,x)\frac{\partial^2 u}{\partial x^2} - r(t,x)u + f(t,x) = 0
\end{equation}
且终端条件 $u(T,x) = g(x)$，则：
\begin{equation}
  u(t,x) = \mathbb{E}\left[\int_t^T e^{-\int_t^s r(\tau,X_\tau)d\tau} f(s,X_s)ds + e^{-\int_t^T r(\tau,X_\tau)d\tau} g(X_T) \mid X_t = x\right]
\end{equation}
\end{theorem}

这建立了 BS PDE 与风险中性期望的等价性——数学上的同一座桥，从两个方向架设。

\section{小结}

随机过程是量化金融的数学心脏。布朗运动提供了基准随机性，伊藤积分使我们能够在"处处不可微"的世界中做微积分，L\'evy 过程将跳跃风险纳入模型，Girsanov 和 Feynman-Kac 构建了 PDE、SDE 和鞅定价的统一框架。这是通向 Q Quant 的必修数学通道。
